\documentclass[aspectratio=169,11pt,svgnames]{beamer}

\usepackage[czech]{babel}
\usepackage[czech=quotes]{csquotes}
\usepackage{graphicx}
\usepackage{enumitem}
\usepackage{amsmath}
\usepackage{mathtools}
\usepackage{float}
\usepackage{tikz}

% Flowchart stuff
\usetikzlibrary{shapes.geometric, arrows, calc}
\tikzstyle{startstop} = [rectangle, rounded corners, minimum width=3cm, minimum
height=1cm,text centered, draw=black, fill=Red!30]
\tikzstyle{io} = [trapezium, trapezium left angle=70, trapezium right angle=110,
minimum width=3cm, minimum height=1cm, text centered, draw=black, fill=Blue!30]
\tikzstyle{process} = [rectangle, minimum width=3cm, minimum height=1cm, text
centered, draw=black, fill=Orange!30]
\tikzstyle{decision} = [diamond, aspect=2, minimum width=3cm, minimum height=.5cm, text
centered, draw=black, fill=Green!30]
\tikzstyle{connector} = [draw, -latex']

\usepackage{pgfopts}
\usepackage{xcolor}
\usepackage{tcolorbox}

\usetheme[
 titlestyle=style2,
 titleformat=smallcaps,
 sectionstyle=plain,
 slidestyle=cyber,
 headingcolor=theme,
 block=transparent
]{trigon}

\title{Algoritmus}
\date{\today}
\author{Adam Klepáč}
\institute[GEVO]{Gymnázium Evolution Jižní Město}
\biglogo[width=.2\textwidth]{logo}
\smalllogo[width=.1\textwidth]{logo}
\titlegraphic{\includegraphics[height=\paperheight]{title.jpg}}

\def\subsectionname{}

% enumerate global settings
\setlist[enumerate,1]{label=\arabic*.}
\setlist[enumerate,2]{label=\alph*)}

% custom colors %
\definecolor{SomeRandomRed}{HTML}{B22028}
\colorlet{tPrim}{LightSeaGreen}

\tcbset{
 boxsep=7pt,
 fonttitle=\sc,
 colframe=tGreyBg,
 colframe=tTheme,
 boxrule=1pt
}

\begin{document}
\titleframe

\begin{frame}
 \frametitle{Obsah}
 \tableofcontents
\end{frame}

\section{Co je algoritmus?}
\subsection{Příklady}

\begin{frame}
 \subsectionpage
\end{frame}

\begin{frame}
 \frametitle{Recept na mramorovou bábovku}
 \begin{block}{Postup}
  \begin{enumerate}
   \item Smíchejte cukr s vejci, přidejte olej, mouku, mléko, prášek do pečiva,
    vanilkový cukr a dobře promíchejte.
   \item Rozdělte těsto do dvou misek. Do jedné misky přidejte podle chuti
    nastrouhanou kůru z citronu.
   \item Do další misky s těstem dejte holandské kakao také podle chuti.
   \item Nalijte do vymazané formy a pečte na 200 °C 45–50 min.
   \item Po upečení ještě horkou vyklopte na talíř. Krájejte vlažnou.
  \end{enumerate}
 \end{block}
\end{frame}

\begin{frame}
 \frametitle{Jak si zavázat tkaničky}
 \begin{block}{Postup}
  \begin{enumerate}
   \item Udělejte uzel.
   \item Udělejte mašličku na pravé tkaničce.
   \item Levou tkaničkou mašličku omotejte.
   \item Pravým ukazováčkem prostrčte pravou mašličku očkem směrem za pravým
    palcem.
   \item Chytněte vrcholy mašliček.
   \item Utáhněte.
  \end{enumerate}
 \end{block}
\end{frame}

\begin{frame}
 \frametitle{Sčítání pod sebou}
 \begin{block}{Postup}<1->
  \begin{enumerate}
   \item<2->\label{item:sum-1} Sečtěte číslice na poslední pozici obou čísel.
   \item<3-> K součtu přičtěte přebytek.
   \item<4-> Výsledný součet napište.
   \item<5-> 
    \begin{enumerate}
     \item Je-li součet větší než 10, nastavte přebytek na 1.
     \item Je-li součet menší než 10, nastavte přebytek na 0.
    \end{enumerate}
   \item<6-> Zapomeňte/smažte poslední číslice obou čísel.
   \item<7-> Pokud ještě oběma číslům zbývají nějaké číslice, opakujte krok
    \ref{item:sum-1}
  \end{enumerate}
 \end{block}
\end{frame}

\subsection{Protipříklady}

\begin{frame}
 \subsectionpage
\end{frame}

\begin{frame}
 \frametitle{Já věřím, že poletím!}
 \begin{block}{Postup}
  \begin{enumerate}
   \item Vystoupejte na nejvyšší vrchol v okolí 5 km.
   \item Vzneste se.
   \item Udělejte si selfie.
   \item Uploadněte je na Instagram.
   \item Leťte 10 km na západ rychlostí zvuku ve vodě.
   \item Přistaňte.
   \item Zkontrolujte liky a commenty.
  \end{enumerate}
 \end{block}
\end{frame}

\begin{frame}
 \frametitle{Nerozhodná navigace}
 \begin{block}{Postup}
  \begin{enumerate}
   \item Jsme připraveni. Řiďte bezpečně.
   \item Pokračujte rovně deset minut až po odbočku Turnov/Liberec/Mnichovo
    Hradiště.
   \item Sjeďte vpravo.
   \item Z kruhového objezdu vyjeďte kterýmkoli výjezdem.
   \item Pokračujte rovně pět minut.
   \item Cíl.
  \end{enumerate}
 \end{block}
\end{frame}

\begin{frame}
 \frametitle{Největší společný násobek}
 \begin{block}{Postup}
  \begin{enumerate}
   \item Máte dána přirozená čísla $A$ a $B$.
   \item Položte $N \coloneqq AB$.
   \item Dokud $A$ dělí $N$ a $B$ dělí $N$, opakujte následující kroky:
    \begin{enumerate}[label=\roman*.]
     \item Položte $N \coloneqq NA$.
     \item Položte $N \coloneqq NB$.
    \end{enumerate}
   \item Vypište $N$.
  \end{enumerate}
 \end{block}
\end{frame}

\subsection{Neformální popis}
\begin{frame}
 \subsectionpage
\end{frame}

\begin{frame}
 \frametitle{Popis pomocí vlastností}
 Algoritmem nazveme \alert{chronologickou sadu příkazů} splňující
 následující tři vlastnosti.
 \pause
 \begin{tcolorbox}[title=Vlastnost 1: splnitelnost,width=.8\textwidth,center]
  Příkazy musejí být splnitelné \alert{z pohledu plnitele}.
 \end{tcolorbox}
\end{frame}

\begin{frame}
 \frametitle{Popis pomocí vlastností}
 Algoritmem nazveme \alert{chronologickou sadu příkazů} splňující
 následující tři vlastnosti.
 \begin{tcolorbox}[title=Vlastnost 2: jednoznačnost,width=.8\textwidth,center]
  Příkazy musejí být jednoznačné \alert{z pohledu plnitele}.
 \end{tcolorbox}
\end{frame}

\begin{frame}
 \frametitle{Popis pomocí vlastností}
 Algoritmem nazveme \alert{chronologickou sadu příkazů} splňující
 následující tři vlastnosti.
 \begin{tcolorbox}[title=Vlastnost 3: konečnost,width=.8\textwidth,center]
  Příkazy musejí vyžadovat jen konečně mnoho prostoru a času.
 \end{tcolorbox}
\end{frame}

\section{Flowchart (vývojový diagram)}

\subsection{Popis}

\begin{frame}
 \subsectionpage
\end{frame}

\begin{frame}
 \frametitle{Části algoritmu}
 \begin{enumerate}
  \item<2-> Začátek/konec (start/end). Ve flowchartu obvykle
   \tikz[baseline=-0.5ex]{\draw[rounded corners] (-0.9,-0.4) rectangle (0.9,0.4)
   {};}
  \item<3-> Vstup/výstup (input/output). Ve flowchartu obvykle 
   \tikz[baseline=-0.5ex,xslant=0.5]{\draw (-0.9,-0.4) rectangle (0.9,0.4) {};}
  \item<4-> Úkon (process). Ve flowchartu obvykle 
   \tikz[baseline=-0.5ex]{\draw (-0.9,-0.4) rectangle (0.9,0.4) {};}
  \item<5-> Rozhodnutí (decision). Ve flowchartu obvykle
   \tikz[baseline=-0.5ex]{\node[draw,diamond,minimum width=1.8cm,minimum
   height=0.8cm] {};}
 \end{enumerate}
\end{frame}

\begin{frame}
 \frametitle{Co je f\mbox{}lowchart?}
 \begin{tcolorbox}[title=Flowchart,width=.8\textwidth,center]
  Flowchart je znázornění průběhu algoritmu jako jeho částí spojených šipkami.
 \end{tcolorbox}
\end{frame}

\subsection{Příklady}
\begin{frame}
 \subsectionpage
\end{frame}

\begin{frame}
 \frametitle{Přihlášení uživatele}
 \begin{figure}[H]
  \centering
  \begin{tikzpicture}[node distance=2cm, scale=0.5, every node/.style={transform
   shape}]
   \node[startstop] (start) at (0,0) {zobrazení přihlašovacího okna};
   \node[io, below of=start] (creds) {uživatelské jméno, heslo};
   \node[decision, below of=creds, yshift=-.5cm] (userexists) {Existuje uživatel?};
   \node[process, right of=userexists,xshift=5cm] (usernotexists) {Oznam
    \enquote{Neplatné uživatelské jméno}.};
   \node[decision, below of=creds, yshift=-3.5cm] (passok) {Správné heslo?};
   \node[process, right of=passok, xshift=5cm] (passnotok) {Oznam
    \enquote{Chybné heslo}.};
   \node[process, below of=passok,yshift=-.5cm] (loguser) {Přihlaš uživatele.};
   \node[startstop,below of=loguser] (nextpage) {jiná stránka};
  
   \path[connector] (start) -- (creds);
   \path[connector] (creds) -- (userexists);
   \path[connector] (userexists) -- node[anchor=south] {ne}
   (usernotexists);
   \path[connector] (userexists) -- node[anchor=east] {ano} (passok);
   \path[connector] (passok) -- node[anchor=south] {ne} (passnotok);
   \path[connector] (passok) -- node[anchor=east] {ano} (loguser);
   \path[connector] (loguser) -- (nextpage);

   \draw[latex'-] (start) -| (usernotexists);
   \draw (passnotok.east) to ++(2,0);
   \draw ($(passnotok.east) + (2,0)) to (11,0);
   \draw (11,0) -- (7,0);
  \end{tikzpicture}
 \end{figure}
\end{frame}

\begin{frame}
 \frametitle{Sčítání pod sebou}
 \begin{figure}[H]
  \centering
  \begin{tikzpicture}[node distance=2cm, scale=0.5, every node/.style={transform
    shape}]
    \node[startstop] (start) at (0,0) {start};
    \node[io,below of=start] (input) {dvě stejně dlouhá čísla};
    \node[decision,below of=input,yshift=-1cm] (digitexists) {Zbývají ještě nějaké číslice?};
    \node[process,right of= start,xshift=7cm] (sum) {Sečti poslední číslice zprava.};
    \node[decision,below of=sum,yshift=-1cm] (greaterthan10) {Je výsledek
     větší než 10?};
    \node[process,below left of=greaterthan10,yshift=-1.5cm,xshift=-1cm] (no) {Nastav přebytek na $0$.};
    \node[process,below right of=greaterthan10,yshift=-1.5cm,xshift=1cm] (yes) {Nastav přebytek na $1$.};
    \node[process,below of=yes,xshift=-2cm] (lastdigit) {Zapiš poslední číslici
     výsledku.};
    \node[process,below of=digitexists,yshift=-1cm] (carry) {Dopiš přebytek.};
    \node[io,below of=carry] (finalsum) {výsledný součet};
    
    \node[startstop,below of=finalsum] (stop) {konec};

    \path[connector] (start) -- (input);
    \path[connector] (input) -- (digitexists);
    \draw (digitexists.east) to ++(2,0);
    \draw ($(digitexists.east) + (2,0)$) to node[midway,xshift=-4mm] {ano}
     (5,0);
    \draw[latex'-] (sum.west) to (5,0);
    \path[connector] (sum) -- (greaterthan10);
    \path[connector] (greaterthan10) -- node[anchor=east] {ne} (no);
    \path[connector] (greaterthan10) -- node[anchor=west] {ano} (yes);
    \path[connector] (no) -- (lastdigit);
    \path[connector] (yes) -- (lastdigit);
    \path[connector] (lastdigit.west) -- (digitexists);
    \path[connector] (digitexists) -- node[anchor=east] {ne} (carry);
    \path[connector] (carry) -- (finalsum);
    \path[connector] (finalsum) -- (stop);
  \end{tikzpicture}
 \end{figure}
\end{frame}

\begin{frame}
 \vfill
 \centering\Huge\sc Díky za pozornost.
 \vfill
\end{frame}

\end{document}
